\section{Experiments}

\subsection{Benchmarks and Setup}
\label{sec:experiments-setup}

We evaluate on seven tool-use benchmark settings that span retrieval, question answering, application control, and workflow automation.
We choose these benchmarks because they stress the cross-block dependency that GraphPTC targets along complementary axes, from short to long horizons, and each carries an executable, rule-based, or judge-based evaluation against a reference.
\textbf{BrowseComp-Plus}~\citep{chen2025browsecompplus} measures evidence-seeking browsing over a fixed corpus and reports answer accuracy.
\textbf{AppWorld}~\citep{trivedi2024appworld} is an executable environment of everyday apps, and we use its normal and challenge splits scored by Task Goal Completion (TGC) and Scenario Goal Completion (SGC).
\textbf{Agent-Diff}~\citep{pysklo2026agentdiff} measures stateful multi-step tasks with an assertion-based pass rate.
\textbf{FanOutQA}~\citep{zhu2024fanoutqa} measures multi-hop, multi-document question answering with a judge accuracy.
\textbf{FRAMES}~\citep{krishna2024frames} measures fact retrieval and reasoning with a judge accuracy.
\textbf{AutomationBench}~\citep{shepard2026automationbench} measures business workflow automation with a pass rate over end-to-end workflows.
The two AppWorld splits give seven evaluation settings in total, and we run every setting with all five backbones.

We instantiate three methods under one shared protocol.
\emph{Direct Tool Calling} uses native per-call function calling, \emph{PTC} uses programmatic tool calling with a persistent runtime and no graph, and \emph{GraphPTC} adds the dynamic execution graph of Section~\ref{sec:method}.
All three share the same tools, prompts, and budgets on each benchmark, and the opportunity rules and feedback budget of GraphPTC are held fixed across all benchmarks and backbones, so a gain reflects the graph and not per-setting tuning.
The five backbones, GPT-5.6-Sol, Claude-Opus-5, DeepSeek-V4-Pro, Kimi-K3, and Qwen3.8-Max, cover closed and open, dense and mixture-of-experts models.

\begin{table}[t]
\centering\scriptsize
\setlength{\abovecaptionskip}{0pt}
\setlength{\belowcaptionskip}{0pt}
\setlength{\tabcolsep}{4pt}
\renewcommand{\arraystretch}{0.95}
\newcommand{\hd}[2]{\begin{tabular}[b]{@{}c@{}}\textbf{#1}\\(#2)\end{tabular}}
\newcommand{\hdd}[3]{\begin{tabular}[b]{@{}c@{}}\textbf{#1}\\\textbf{#2}\\(#3)\end{tabular}}
\caption{Main results on seven tool-use benchmark settings across five backbones. Each column reports the primary metric of one setting, with answer accuracy for BrowseComp-Plus, Task Goal Completion (TGC) for the two AppWorld splits, pass rate for Agent-Diff and AutomationBench, and judge accuracy for FanOutQA and FRAMES. \textbf{Bold} marks the best method per backbone and benchmark. GraphPTC improves over both Direct Tool Calling and PTC on every benchmark.}
\label{tab:main}
\resizebox{0.93\linewidth}{!}{%
\begin{tabular}{l ccccccc}
\toprule
\textbf{Method}
 & \hdd{BrowseComp-}{Plus}{Acc}
 & \hdd{AppWorld-}{Normal}{TGC}
 & \hdd{AppWorld-}{Challenge}{TGC}
 & \hd{Agent-Diff}{Pass}
 & \hd{FanOutQA}{Acc}
 & \hd{FRAMES}{Acc}
 & \hdd{Automation-}{Bench}{Pass} \\
\midrule
\multicolumn{8}{c}{\textit{\textbf{GPT-5.6-Sol}}} \\
Direct Tool Calling & 73.49 & 82.10 & 81.80 & 51.34 & 55.81 & 77.31 & 20.83 \\
PTC                 & 80.00 & 89.90 & 90.60 & 39.73 & 52.58 & 73.42 & 29.17 \\
\textbf{GraphPTC}   & \best{80.60} & \best{91.70} & \best{93.00} & \best{51.49} & \best{60.00} & \best{79.00} & \best{31.00} \\
\midrule
\multicolumn{8}{c}{\textit{\textbf{Claude-Opus-5}}} \\
Direct Tool Calling & 43.98 & 94.00 & 80.80 & 75.14 & 77.74 & 84.95 & 28.67 \\
PTC                 & 10.96 & 59.50 & 79.10 & 68.06 & 76.77 & 85.19 & 32.50 \\
\textbf{GraphPTC}   & \best{49.04} & \best{95.80} & \best{95.00} & \best{79.17} & \best{78.71} & \best{85.56} & \best{33.50} \\
\midrule
\multicolumn{8}{c}{\textit{\textbf{DeepSeek-V4-Pro}}} \\
Direct Tool Calling & 4.58  & 81.00 & 65.70 & 70.09 & 57.42 & 81.43 & 21.50 \\
PTC                 & 3.49  & 88.70 & 80.10 & 72.47 & 58.06 & 80.83 & 17.00 \\
\textbf{GraphPTC}   & \best{13.37} & \best{89.90} & \best{81.50} & \best{73.96} & \best{62.26} & \best{81.55} & \best{22.83} \\
\midrule
\multicolumn{8}{c}{\textit{\textbf{Kimi-K3}}} \\
Direct Tool Calling & 8.92  & 72.00 & 46.50 & 77.53 & 49.03 & 73.54 & 27.50 \\
PTC                 & 2.89  & 85.70 & 83.20 & 78.12 & 28.71 & 75.85 & 28.00 \\
\textbf{GraphPTC}   & \best{10.48} & \best{86.30} & \best{84.90} & \best{78.72} & \best{50.65} & \best{76.46} & \best{28.83} \\
\midrule
\multicolumn{8}{c}{\textit{\textbf{Qwen3.8-Max}}} \\
Direct Tool Calling & 34.46 & 82.70 & 71.70 & 74.26 & 65.81 & 79.98 & 31.67 \\
PTC                 & 39.76 & 95.20 & 85.40 & 78.12 & 59.68 & 80.10 & 32.38 \\
\textbf{GraphPTC}   & \best{41.93} & \best{97.00} & \best{87.50} & \best{80.06} & \best{66.77} & \best{82.40} & \best{42.41} \\
\bottomrule
\end{tabular}%
}
\end{table}


\subsection{Main Results}

\textbf{Capability on Application Tool Calling.}\quad
We evaluate \method on the AppWorld normal and challenge splits to assess its \emph{multi-step application control}, where a task drives many app calls whose write effects must persist and compose into a final goal.
As shown in Table~\ref{tab:main}, \method improves over both Direct Tool Calling and PTC on both splits for every backbone, and its five-backbone average gain over Direct Tool Calling grows from $9.8$ points on the normal split to $19.1$ points on the challenge split.
We attribute this pattern to modeling state as typed nodes, so a later block links to an earlier write instead of re-deriving it from text, and Section~\ref{sec:analysis} tests this attribution directly.

\textbf{Capability on Multi-Hop Tool Calling.}\quad
We evaluate \method on BrowseComp-Plus, FanOutQA, and FRAMES to assess its \emph{multi-hop retrieval and reasoning}, where an answer depends on composing evidence gathered across many calls.
As shown in Table~\ref{tab:main}, \method improves over both baselines on all three benchmarks for every backbone, and it lifts the collapsed PTC accuracy of Claude-Opus-5 on BrowseComp-Plus from $10.96$ to $49.04$, above the $43.98$ of Direct Tool Calling.
This recovery reflects cross-block dependency linking, which marks equivalent results and lets the agent stop re-issuing queries whose content it already holds.

\textbf{Capability on Long-Horizon Tool Calling.}\quad
We evaluate \method on Agent-Diff and AutomationBench to assess its \emph{stateful workflow automation}, where a task completes only when a chain of dependent write operations all take effect.
As shown in Table~\ref{tab:main}, \method improves over both baselines on both benchmarks for every backbone, and on AutomationBench it turns the DeepSeek-V4-Pro case where PTC loses $4.5$ points to Direct Tool Calling into a gain of $1.3$ points.
Here progress assessment separates a genuine state change from a re-executed no-op, so the agent advances the workflow instead of restating a completed step.

Across the panel, \method is the best method on every one of the seven settings for all five backbones, improving over PTC by $6.3$ points and over Direct Tool Calling by $6.8$ points on average.
PTC itself falls below Direct Tool Calling $14$ times over the same settings and backbones, losing $34.5$ points on Claude-Opus-5 in the AppWorld normal split, so a programmatic runtime needs structure to pay off and is not a uniform upgrade.
